\section{Verification methodology}
\label{sec:verification}

\subsection{What one verification pass costs}

The model is verified by a suite of scenarios, each of which constructs a
world, drives the simulation, and asserts on the outcome. At the revision
described here the suite comprises \textbf{130 scenarios} making \textbf{594
assertions}. Executing it performs \textbf{12{,}322 simulated days} --- about
34 simulated years --- and generates \textbf{265 independent worlds}, each a
$128\times128$ field with a full solar raycast.

These figures are counted, not estimated: \texttt{e1-model-scale.mjs} wraps the
simulation's tick and world-generation entry points and runs the real suite.
We report them because they are the honest unit of reproducibility for a model
of this kind. A statement that the model "passes its tests" means little; a
statement that it survives twelve thousand simulated days across two hundred
and sixty-five independently generated worlds is a claim a reader can weigh.

\subsection{Scenarios as executable claims}

Each scenario carries a machine-readable statement of \emph{why} it exists
alongside its assertions. This is not documentation but a design constraint: a
scenario whose rationale cannot be written down is usually asserting an
implementation detail rather than a property of the model, and several were
deleted on those grounds during development.

\subsection{On counting simulation runs}

A reader may reasonably ask how many times the simulation was run in the course
of building the model. We can answer that in three parts, and only two of them
honestly.

\begin{enumerate}
\item \textbf{The cost of one pass} is exact and reproducible, and is given
  above.
\item \textbf{Runs from the point of instrumentation onward} are recorded.
  Every experiment appends to \texttt{results/ledger.json} a record of what it
  executed --- simulated days, worlds generated, and the source revision --- so
  the cost of reproducing any figure in this paper is itself a committed
  artefact.
\item \textbf{Runs before instrumentation are not recoverable.} The model was
  developed over roughly six weeks and 31 commits touching its directory,
  during which the suite was executed many times, but no count was kept and we
  decline to estimate one. Reconstructing it from commit timestamps would
  produce a number with the shape of evidence and none of the substance.
\end{enumerate}

We state this plainly because provenance is exactly the kind of claim that
invites a plausible fabrication, and because the fix --- instrument first --- is
cheap and was applied here only in retrospect. The ledger closes the gap
prospectively, not retrospectively.
